
% ==============================================================================
% CAPÍTULO 2
% LOCALIZACIÓN ACÚSTICA: MODELOS GEOMÉTRICOS Y ESTIMACIÓN CINEMÁTICA
% ==============================================================================

\clearpage
\section{LOCALIZACIÓN ACÚSTICA: MODELOS GEOMÉTRICOS Y ESTIMACIÓN CINEMÁTICA}
\label{chap:modelos_inversion}

% ==============================================================================
% 2.1. PLANTEAMIENTO DEL PROBLEMA INVERSO
% ==============================================================================

\subsection{El problema inverso de la localización acústica}

La localización acústica consiste en estimar propiedades geométricas o
cinemáticas de una fuente sonora a partir de las señales registradas por uno
o varios receptores. El problema directo describe la propagación de la onda
desde la fuente hasta los sensores; el problema inverso utiliza las señales
medidas para inferir magnitudes como la distancia, la dirección de llegada
y la velocidad de la fuente.

La información disponible depende del observable físico utilizado. Entre
los métodos habituales se encuentran:

\begin{itemize}
    \item \textbf{Métodos basados en amplitud o potencia recibida (RSS):}
    relacionan la amplitud de la señal o una magnitud proporcional a la
    potencia recibida con la distancia entre la fuente y el receptor.
    Su validez depende de la directividad de la fuente, la respuesta de
    los transductores y las condiciones de propagación
    \cite{kinsler2000fundamentals,zhuang2016evaluation}.

    \item \textbf{Métodos basados en tiempo de llegada (ToA):}
    utilizan el tiempo de propagación entre la fuente y el receptor.
    Permiten estimar una distancia absoluta cuando se conoce el instante
    de emisión y se han caracterizado los retardos instrumentales.

    \item \textbf{Métodos basados en diferencias de tiempo de llegada
    (TDOA):}
    utilizan la diferencia entre los tiempos registrados por distintos
    receptores. La magnitud obtenida corresponde a una diferencia de
    distancias y no, por sí sola, a una posición única
    \cite{knapp1976generalized,brandstein2001microphone}.

    \item \textbf{Métodos basados en diferencias de fase (PDoA):}
    relacionan la diferencia de fase de una señal armónica recibida por
    dos sensores con la diferencia de camino de propagación. Su precisión
    depende de la frecuencia, la separación entre sensores, la coherencia
    entre canales y la ambigüedad de fase
    \cite{brandstein2001microphone,rautenberg2018phase}.

    \item \textbf{Métodos basados en el efecto Doppler:}
    utilizan el desplazamiento de frecuencia producido por el movimiento
    relativo entre la fuente y los receptores. Bajo una geometría conocida,
    las frecuencias recibidas pueden relacionarse con la velocidad de la
    fuente \cite{kinsler2000fundamentals,chestnut1982estimation}.
\end{itemize}

En WALO se consideran cinco observables principales: tiempo de llegada
(ToA), diferencia de tiempos de llegada (TDOA), diferencia de fase
intercanal (PDoA), frecuencia recibida y diferencia de frecuencia asociada
al efecto Doppler. También se obtiene una magnitud digital relacionada
con la amplitud cuadrática de la señal recibida.

Estos observables no contienen la misma información ni requieren las
mismas condiciones para ser interpretados. Por ello, cada modelo de
inversión se formula por separado y se especifican las hipótesis necesarias
para relacionar la magnitud medida con una propiedad de la fuente.

El propósito de este capítulo es establecer esas relaciones físicas.
La caracterización experimental de los estimadores y sus limitaciones
se presenta posteriormente.

% ==============================================================================
% 2.2. ESTIMACIÓN DE DISTANCIA
% ==============================================================================

\subsection{Estimación de distancia}

La distancia radial $r$ entre una fuente emisora y un receptor puede
estimarse mediante el tiempo de propagación o a partir de la atenuación
de la señal recibida. Ambos procedimientos utilizan información física
diferente y presentan fuentes de error distintas. No obstante, comparten
factores como la temperatura del medio, la geometría de propagación y
la respuesta electroacústica del sistema.

\begin{figure}[!htbp]
    \centering
    % figures/fig4_geometria_to_a.tex
\begin{tikzpicture}[
    >=Stealth,
    stand/.style={
        line width=1.5pt, draw=black!75
    },
    baseplate/.style={
        draw=black!80, fill=black!25, rounded corners=1pt
    },
    sensor/.style={
        circle, draw=black!85, fill=black!12, line width=1.1pt,
        minimum size=0.60cm, inner sep=0pt
    },
    % Fuente en naranja unificado
    source/.style={
        circle, draw=black!85, fill=orange!80, line width=1.1pt,
        minimum size=0.60cm, inner sep=0pt
    },
    tabletop/.style={
        draw=black!75, line width=0.9pt, pattern=north east lines, pattern color=black!25
    },
    raydir/.style={
        ->, line width=1.2pt, draw=UdeAGreen!75!black
    },
    rayrefl/.style={
        ->, line width=1.1pt, dashed, draw=red!70!black
    }
]

    % =========================================================================
    % 1. SUPERFICIE HORIZONTAL DE LA MESA (y = 0)
    % =========================================================================
    \draw[tabletop] (-5.6, -0.38) rectangle (5.6, 0.0);
    \node[below=4pt, font=\footnotesize\sffamily] at (0, -0.38) {Superficie reflectora horizontal (Mesa / Suelo)};

    % =========================================================================
    % 2. SOPORTES Y TRANSDUCTORES (x = -3.6 y x = +3.6, y = 2.6)
    % =========================================================================
    \coordinate (Mic) at (-3.6, 2.6);       % Receptor R
    \coordinate (Spk) at (3.6, 2.6);        % Emisor S
    \coordinate (Bounce) at (0.0, 0.0);     % Punto de rebote en la mesa

    % Soporte y cápsula Receptor R (Gris)
    \draw[baseplate] (-4.05, 0.0) rectangle (-3.15, 0.14);
    \draw[stand] (-3.6, 0.14) -- (Mic);
    \node[sensor] (NodeMic) at (Mic) {};
    \node[left=8pt, font=\small\sffamily\bfseries] at (Mic) {Receptor ($R$)};

    % Soporte y cápsula Fuente S (Naranja)
    \draw[baseplate] (3.15, 0.0) rectangle (4.05, 0.14);
    \draw[stand] (3.6, 0.14) -- (Spk);
    \node[source] (NodeSpk) at (Spk) {};
    \node[right=8pt, font=\small\sffamily\bfseries] at (Spk) {Fuente ($S$)};

    % Arcos acústicos de radiación en naranja unificado
    \draw[line width=0.85pt, orange!85!black] ($(Spk)+(-0.45, 0.3)$) arc (145:215:0.5);
    \draw[line width=0.85pt, orange!85!black] ($(Spk)+(-0.75, 0.45)$) arc (145:215:0.8);
    \draw[line width=0.85pt, orange!85!black] ($(Spk)+(-1.05, 0.6)$) arc (145:215:1.1);

    % =========================================================================
    % 3. TRAYECTORIAS ONDULATORIAS Y COTA CENTRAL DE ALTURA h
    % =========================================================================
    % Trayectoria directa horizontal
    \draw[raydir] (NodeSpk.west) -- node[above=4pt, font=\small\sffamily] {Onda directa: $r_{\text{dir}} = r$} (NodeMic.east);

    % Trayectoria reflejada especular en V
    \draw[rayrefl] (NodeSpk.south) -- node[below right=2pt, font=\footnotesize\sffamily, text=red!75!black] {$r_{\text{refl}}/2$} (Bounce);
    \draw[rayrefl] (Bounce) -- node[below left=2pt, font=\footnotesize\sffamily, text=red!75!black] {$r_{\text{refl}}/2$} (NodeMic.south);
    \fill[red!70!black] (Bounce) circle (2.0pt);

    % Eje vertical de altura h
    \draw[dashed, draw=black!45, line width=0.6pt] (0.0, 0.0) -- (0.0, 2.6);
    \draw[<->, draw=black!75, line width=0.65pt] (-0.35, 0.05) -- node[left=1pt, font=\small\sffamily] {$h$} (-0.35, 2.55);

    % =========================================================================
    % 4. CAJA METROLÓGICA (y = 5.0)
    % =========================================================================
    \node[draw=black!75, fill=white, rounded corners=4pt, line width=0.8pt, inner sep=8pt] at (0.0, 5.0) {
        \footnotesize
        \begin{tabular}{@{} l l @{}}
            \multicolumn{2}{@{} l @{}}{\textbf{Geometría de Reflexión y Condición Cuasi-Anecoica:}} \\[4pt]
            Trayectoria directa: & $r_{\text{dir}} = r$ \\
            Trayectoria reflejada: & $r_{\text{refl}} = 2\sqrt{(r/2)^2 + h^2} = \sqrt{r^2 + 4h^2}$ \\
            Diferencial de camino: & $\Delta r_{\text{refl}} = \sqrt{r^2 + 4h^2} - r$ \\[4pt]
            \multicolumn{2}{@{} l @{}}{\textcolor{UdeAGreen!80!black}{\textbf{Condición de corte:}} \ $c(T) \cdot T_{\text{gate}} < \Delta r_{\text{refl}} \implies \text{Anulación de eco }(c_{\text{room}} \to 0)$}
        \end{tabular}
    };

\end{tikzpicture}
    \vspace{0.25cm}
    \caption{Geometría de propagación directa y reflexión especular en una
    superficie horizontal. Para una fuente y un receptor situados a la
    misma altura $h$, el método de la fuente imagen permite calcular la
    diferencia entre las longitudes de los trayectos directo y reflejado.}
    \label{fig:geometria_cuasi_anecoica}
\end{figure}

\subsubsection{Estimación temporal: tiempo de llegada (ToA)}

En condiciones de propagación directa, el tiempo de vuelo de una onda
acústica entre una fuente y un receptor separados una distancia $r$ es

\begin{equation}
    t_{\mathrm{vuelo}} = \frac{r}{c(T)},
    \label{eq:tiempo_vuelo}
\end{equation}

donde $c(T)$ es la velocidad del sonido en el medio. Para aire seco,
una aproximación de su dependencia con la temperatura es

\begin{equation}
    c(T) = 331.3
    \sqrt{1+\frac{T}{273.15}},
    \label{eq:velocidad_sonido_temperatura}
\end{equation}

con $T$ expresada en grados Celsius. Esta expresión no incorpora
correcciones por humedad, composición del aire ni gradientes térmicos.

Si el instante de emisión es conocido y el retardo instrumental se
representa mediante $t_{\mathrm{off}}$, la distancia puede estimarse como

\begin{equation}
    r(t) =
    c(T)\left[t_{\mathrm{arr}}(t)-t_{\mathrm{em}}(t)
    -t_{\mathrm{off}}\right],
    \label{eq:distancia_toa_absoluta}
\end{equation}

donde $t_{\mathrm{arr}}$ es el instante de llegada detectado,
$t_{\mathrm{em}}$ es el instante de emisión y $t_{\mathrm{off}}$
agrupa los retardos sistemáticos de la cadena electroacústica y de
adquisición.

Una calibración a distancia conocida $r_{\mathrm{ref}}$ permite
determinar un valor efectivo del retardo:

\begin{equation}
    t_{\mathrm{off}} =
    t_{\mathrm{arr}}^{(\mathrm{ref})}
    -t_{\mathrm{em}}^{(\mathrm{ref})}
    -\frac{r_{\mathrm{ref}}}{c(T)}.
    \label{eq:calibracion_toffset}
\end{equation}

Esta calibración corrige el retardo global bajo las condiciones del
ensayo de referencia. No separa necesariamente las contribuciones
individuales ni garantiza que el retardo permanezca constante ante
cambios de frecuencia, temperatura, amplitud o configuración del
procesamiento. Su estabilidad debe comprobarse cuando cambian las
condiciones de medida.

La conversión de una variación temporal en una variación de distancia
se expresa como

\begin{equation}
    \delta r = c(T)\,\delta t.
    \label{eq:conversion_resolucion_toa}
\end{equation}

A modo de referencia, una variación temporal de
$\SI{1}{\micro\second}$ corresponde aproximadamente a
$\SI{0.34}{\milli\metre}$ de recorrido acústico para
$c=\SI{343.21}{\metre\per\second}$. Este resultado expresa una
equivalencia cinemática entre tiempo y distancia; no constituye por sí
mismo la incertidumbre ni la exactitud de la estimación ToA. Estas
últimas dependen del ruido, la forma del pulso, el criterio de detección,
la frecuencia de muestreo, la interpolación y la calibración del retardo.

Es importante distinguir el ToA del TDOA. El primero requiere una
referencia temporal asociada a la emisión para obtener una distancia
absoluta. El segundo utiliza diferencias entre receptores y permite
estimar diferencias de distancias, incluso cuando no se conoce el
instante absoluto de emisión.

\subsubsection{Estimación mediante amplitud y energía recibidas}

En un medio homogéneo y sin pérdidas, una fuente puntual que radia
aproximadamente de forma esférica produce una amplitud de presión
acústica proporcional a $1/r$ en campo lejano. Esta dependencia se
relaciona con la conservación de la potencia radiada sobre superficies
esféricas cuya área crece como $4\pi r^2$
\cite{kinsler2000fundamentals,morse1968theoretical}.

Si la respuesta del sistema de emisión y recepción se mantiene
constante, una magnitud eléctrica proporcional a la amplitud acústica
puede modelarse empíricamente como

\begin{equation}
    A(r) = \frac{k_A}{r},
    \label{eq:modelo_amplitud_radial}
\end{equation}

de donde se obtiene la estimación

\begin{equation}
    r_{\mathrm{amp}} =
    \frac{k_A}{A(r)}.
    \label{eq:inversion_amplitud}
\end{equation}

La constante $k_A$ debe determinarse mediante calibración y recoge
la intensidad de emisión, el acoplamiento electroacústico y la ganancia
del canal de recepción. Por tanto, $A(r)$ debe entenderse como una
amplitud eléctrica medida, salvo que exista una calibración adicional
que permita expresarla en unidades de presión acústica.

De manera análoga, si una magnitud digital acumulada es proporcional
al cuadrado de la amplitud durante una ventana temporal de duración
fija, puede utilizarse el modelo empírico

\begin{equation}
    \mathcal{E}_{\mathrm{dig}}(r) =
    \frac{k_E}{r^2},
    \label{eq:modelo_energia_radial}
\end{equation}

y su inversión

\begin{equation}
    r_{\mathrm{energy}} =
    \sqrt{\frac{k_E}{\mathcal{E}_{\mathrm{dig}}(r)}}.
    \label{eq:inversion_energia}
\end{equation}

Aquí $\mathcal{E}_{\mathrm{dig}}$ representa una magnitud digital
relacionada con la amplitud cuadrática de las muestras. Su interpretación
cuantitativa depende de la definición, normalización y ventana temporal
empleadas en la implementación. No debe interpretarse directamente
como energía acústica absoluta en julios. La constante $k_E$ se determina
experimentalmente y depende de la definición de la magnitud digital,
de la ventana de integración y de la respuesta del sistema.

Las relaciones anteriores requieren condiciones de propagación
próximas a las de campo libre, una fuente con emisión suficientemente
estable, una respuesta de recepción reproducible y una ventana temporal
comparable entre mediciones. La atenuación atmosférica, la directividad
de la fuente, las reflexiones y las variaciones de ganancia pueden
apartar las observaciones de las leyes ideales $1/r$ y $1/r^2$.

Por ello, estas expresiones deben utilizarse como modelos de calibración
dentro del intervalo de distancias caracterizado, no como leyes
universales para cualquier recinto. La compatibilidad de los datos
experimentales con estas dependencias tampoco demuestra por sí sola
que la fuente sea isotrópica o que el recinto sea anecoico.

\subsubsection{Reflexión especular y selección de la ventana temporal}

Las reflexiones de las superficies próximas pueden superponerse a la
señal directa y modificar su amplitud y fase. Para estudiar una reflexión
especular en una superficie horizontal, se considera una fuente y un
receptor situados a la misma altura $h$ sobre dicha superficie y
separados horizontalmente una distancia $r$.

El trayecto directo tiene longitud

\begin{equation}
    r_{\mathrm{dir}} = r.
\end{equation}

Mediante el método de la fuente imagen, el trayecto reflejado se
representa como una línea recta entre el receptor y la imagen de la
fuente situada a una distancia vertical $2h$ respecto de la fuente real.
Su longitud es

\begin{equation}
    r_{\mathrm{refl}} = \sqrt{r^2+4h^2}.
\end{equation}

La diferencia entre ambos recorridos resulta

\begin{equation}
    \Delta r_{\mathrm{refl}} =
    \sqrt{r^2+4h^2}-r,
    \label{eq:delta_r_refl}
\end{equation}

y el retraso de llegada de esa reflexión respecto a la onda directa es

\begin{equation}
    \Delta t_{\mathrm{refl}} =
    \frac{\Delta r_{\mathrm{refl}}}{c(T)}.
    \label{eq:delta_t_refl}
\end{equation}

Si la ventana de integración comienza con la llegada de la señal
directa, una condición geométrica para que termine antes de la llegada
de esa reflexión es

\begin{equation}
    T_{\mathrm{gate}} <
    \frac{\sqrt{r^2+4h^2}-r}{c(T)}.
    \label{eq:criterio_corte_anecoico}
\end{equation}

Esta condición solo considera la reflexión especular de la superficie
horizontal descrita. No excluye reflexiones procedentes de otras
superficies ni garantiza que la respuesta del recinto sea anecoica.
Además, en una adquisición real deben considerarse la duración del
pulso, el tiempo de establecimiento del transductor y la posición
efectiva de inicio de la ventana.

En consecuencia, el criterio constituye una condición geométrica de
selección temporal para una reflexión particular. Cuando la ventana
de observación queda contenida entre la llegada directa y la primera
reflexión relevante, puede hablarse de condiciones de propagación
cuasi libres de reflexiones durante esa ventana, sin atribuir esa
propiedad al recinto completo.

% ==============================================================================
% 2.3. ESTIMACIÓN ANGULAR
% ==============================================================================

\subsection{Estimación angular mediante diferencias intercanal}

La estimación angular requiere relacionar las señales registradas por
dos receptores separados una distancia conocida $d$. Se considera un
arreglo lineal situado sobre el eje $x$, con centro en el origen,
mientras que el eje $y$ define la dirección normal al arreglo. El
ángulo $\theta$ se mide respecto a dicho eje normal, de modo que
$\theta=0$ corresponde a la dirección broadside.

\begin{figure}[!htbp]
    \centering
    % figures/fig5_geometria_2d_angular.tex
\begin{tikzpicture}[
    >=Stealth,
    sensor/.style={
        circle, draw=black!85, fill=black!12, line width=1pt,
        minimum size=0.52cm, inner sep=0pt
    },
    source/.style={
        circle, draw=black!85, fill=orange!80, line width=1.1pt,
        minimum size=0.55cm, inner sep=0pt
    },
    axis/.style={
        ->, line width=0.85pt, draw=black!65
    }
]

    % =========================================================================
    % 1. EJES CARTESIANOS (Origen en x = 0.5)
    % =========================================================================
    \coordinate (Orig) at (0.5, 0.0);

    % Ejes coordenados alargados
    \draw[axis] (-5.6, 0.0) -- (5.6, 0.0) node[right=2pt, font=\small\sffamily] {$x$ (Transversal)};
    \draw[axis] (0.5, -0.6) -- (0.5, 5.7) node[above=2pt, font=\small\sffamily] {$y$ (Broadside)};

    % =========================================================================
    % 2. SENSORES RECEPTORES: R1(-d/2, 0) y R2(+d/2, 0)
    % =========================================================================
    % d/2 = 2.0 => R1 en -1.5; R2 en 2.5
    \coordinate (R1) at (-1.5, 0.0);
    \coordinate (R2) at (2.5, 0.0);

    \node[sensor] (NodeR1) at (R1) {};
    \node[below=6pt, font=\small\sffamily] at (R1) {$R_1\left(-\frac{d}{2}, 0\right)$};

    \node[sensor] (NodeR2) at (R2) {};
    \node[below=6pt, font=\small\sffamily] at (R2) {$R_2\left(+\frac{d}{2}, 0\right)$};

    % Cota de línea de base d (inferior con amplio respiro)
    \draw[<->, draw=black!70, line width=0.65pt] (-1.5, -1.30) -- node[below=2pt, font=\footnotesize\sffamily] {Línea de base $d = \norm{\vec{r}_2 - \vec{r}_1}$} (2.5, -1.30);
    \draw[dotted, line width=0.5pt, draw=black!45] (-1.5, 0.0) -- (-1.5, -1.30);
    \draw[dotted, line width=0.5pt, draw=black!45] (2.5, 0.0) -- (2.5, -1.30);

    % =========================================================================
    % 3. FUENTE EMISORA S(x, y) (Ondas hacia abajo, texto limpio arriba)
    % =========================================================================
    \coordinate (S) at (1.7, 4.4);

    % Proyecciones ortogonales hacia los ejes
    \draw[dotted, line width=0.65pt, draw=black!55] (1.7, 4.4) -- (1.7, 0.0) node[below=3pt, font=\small\sffamily] {$x$};
    \draw[dotted, line width=0.65pt, draw=black!55] (1.7, 4.4) -- (0.5, 4.4) node[left=2pt, font=\small\sffamily] {$y$};

    % Frentes de onda concéntricos radiados HACIA ABAJO (hacia los receptores)
    \draw[line width=0.75pt, orange!85!black] ($(S)+(215:0.45)$) arc (215:285:0.45);
    \draw[line width=0.75pt, orange!85!black] ($(S)+(215:0.80)$) arc (215:285:0.80);
    \draw[line width=0.75pt, orange!85!black] ($(S)+(215:1.15)$) arc (215:285:1.15);

    % Nodo Fuente (el rótulo queda arriba, totalmente libre de ondas)
    \node[source] (NodeS) at (S) {};
    \node[above=6pt, font=\small\sffamily\bfseries] at (S) {$S(x, y)$};

    % =========================================================================
    % 4. VECTORES DE PROPAGACIÓN r1, r2, r Y ÁNGULO THETA
    % =========================================================================
    % Vector r1
    \draw[->, line width=1.15pt, draw=UdeAGreen!75!black] (NodeR1) -- node[pos=0.45, above left=1pt, font=\small] {$\vec{r}_1$} (NodeS);

    % Vector r2
    \draw[->, line width=1.15pt, draw=blue!70!black] (NodeR2) -- node[pos=0.55, right=2pt, font=\small] {$\vec{r}_2$} (NodeS);

    % Vector central r
    \draw[->, dashed, line width=0.85pt, draw=black!70] (Orig) -- node[pos=0.65, below right=1pt, font=\small] {$\vec{r}$} (NodeS);

    % Arco del ángulo de incidencia theta respecto al eje vertical y
    \draw[line width=0.8pt, draw=black!85] ($(Orig) + (0.0, 2.2)$) arc (90:74.8:2.2);
    \node[font=\small\sffamily\bfseries] at ($(Orig) + (82.5:2.55)$) {$\theta$};

    % =========================================================================
    % 5. HIPÉRBOLA TDOA
    % =========================================================================
    \draw[thick, orange!90!black, domain=0.4:5.2, smooth, variable=\y]
        plot ({ 0.5 + 0.4845 * sqrt(1 + (\y/1.940)*(\y/1.940)) }, {\y});
    % Rótulo en la parte superior derecha, libre de ondas y de la fuente
    \node[orange!90!black, font=\footnotesize\sffamily, right=3pt] at (1.92, 5.25) {Hipérbola TDOA ($\Delta r = r_1 - r_2$)};

    % =========================================================================
    % 6. RECUADRO ESQUINADO EN EL EXTREMO SUPERIOR IZQUIERDO
    % =========================================================================
    \node[draw=black!75, fill=white, rounded corners=4pt, line width=0.8pt, inner sep=7pt] at (-3.3, 4.2) {
        \footnotesize
        \begin{tabular}{@{} l @{}}
            \textbf{Inversión Angular (PDoA):} \\[3pt]
            Retardo espacial: \ $\Delta r = d \sin\theta$ \\
            Desfase medido: \ $\Delta\Phi = \frac{2\pi f_0}{c(T)} d \sin\theta$ \\[5pt]
            \textcolor{UdeAGreen!80!black}{\textbf{Cota de no-aliasing espacial:}} \\[2pt]
            $d \le \frac{\lambda_0}{2} = \frac{c(T)}{2 f_0}$
        \end{tabular}
    };

\end{tikzpicture}
    \vspace{0.25cm}
    \caption{Geometría del arreglo lineal de dos receptores. El ángulo
    de incidencia se mide respecto al eje normal al arreglo. La
    diferencia de camino entre los receptores permite formular
    estimadores angulares basados en fase o en tiempos de llegada.}
    \label{fig:geometria_2d_angular}
\end{figure}

\subsubsection{Diferencia de camino y diferencia de tiempos (TDOA)}

En el régimen de campo lejano, las direcciones de propagación hacia
los dos receptores pueden considerarse aproximadamente paralelas.
La diferencia de camino entre los sensores se expresa como

\begin{equation}
    \Delta r = r_1-r_2 = d\sin\theta,
    \label{eq:diferencia_camino_tdoa}
\end{equation}

donde $r_1$ y $r_2$ son las distancias desde la fuente hasta los
receptores $R_1$ y $R_2$, respectivamente. Esta convención de signo
corresponde a $R_1=(-d/2,0)$ y $R_2=(+d/2,0)$.

Como el tiempo de propagación es la distancia dividida por la
velocidad del sonido, la diferencia de tiempos de llegada satisface

\begin{equation}
    \Delta t =
    t_1-t_2 =
    \frac{d}{c(T)}\sin\theta.
    \label{eq:diferencia_tiempos_tdoa}
\end{equation}

Despejando el ángulo se obtiene

\begin{equation}
    \theta_{\mathrm{TDOA}} =
    \arcsin\left(
    \frac{c(T)\Delta t}{d}
    \right).
    \label{eq:inversion_angular_tdoa}
\end{equation}

Esta expresión supone una onda incidente aproximadamente plana,
una geometría conocida y canales cuya diferencia de retardo
instrumental ha sido corregida. La estimación también requiere que
la diferencia temporal sea compatible con el intervalo físico
$[-d/c(T),d/c(T)]$.

El TDOA no debe confundirse con el ToA absoluto. En esta geometría,
la diferencia temporal permite estimar una dirección de llegada,
pero no determina la distancia radial a la fuente. Además, una
única línea base proporciona una estimación angular en el plano
definido por el arreglo, no una posición bidimensional única.

\subsubsection{Diferencia de fase intercanal (PDoA)}

Para una señal armónica de frecuencia $f_0$, la diferencia de fase
debida a la propagación es

\begin{equation}
    \Delta\Phi =
    \Phi_2-\Phi_1 =
    \frac{2\pi f_0}{c(T)}d\sin\theta.
    \label{eq:diferencia_fase_intercanal}
\end{equation}

La relación anterior presupone que se ha corregido la diferencia de
fase instrumental entre canales y que la fase se interpreta con una
convención consistente. La fase medida es periódica módulo $2\pi$;
por tanto, su inversión requiere seleccionar una rama de fase
compatible con la geometría y con el intervalo angular de interés.

Cuando esa rama es conocida y el argumento se encuentra dentro del
dominio físico, la estimación angular es

\begin{equation}
    \theta_{\mathrm{PDoA}} =
    \arcsin\left(
    \frac{c(T)\Delta\Phi}{2\pi f_0d}
    \right).
    \label{eq:inversion_angular_fase}
\end{equation}

En la implementación numérica puede limitarse el argumento del
arcoseno para evitar errores de dominio debidos al redondeo.
Sin embargo, si el valor medido queda fuera del intervalo $[-1,1]$
por una discrepancia física o por ruido, el recorte numérico no
resuelve la incompatibilidad: puede ocultarla y sesgar la estimación
hacia un ángulo extremo. Estos casos deben identificarse como
mediciones no compatibles con el modelo o tratarse mediante un
procedimiento estadístico explícito.

La ecuación también muestra que, para una frecuencia fija, la
diferencia de fase y la diferencia de tiempos están relacionadas por

\begin{equation}
    \Delta\Phi = 2\pi f_0\Delta t.
    \label{eq:relacion_fase_tiempo}
\end{equation}

Aunque ambas magnitudes contienen información sobre la diferencia
de camino, sus estimadores no son equivalentes desde el punto de
vista experimental. El TDOA depende de la identificación temporal
de los arribos, mientras que el PDoA depende de la fase de la
componente armónica y de su coherencia entre canales.

\subsubsection{Separación entre sensores y ambigüedad espacial}

La longitud de onda de una señal armónica es

\begin{equation}
    \lambda_0 = \frac{c(T)}{f_0}.
    \label{eq:longitud_onda}
\end{equation}

La diferencia de camino geométrica de un arreglo lineal satisface
$|\Delta r|\leq d$. Para evitar ambigüedades espaciales de fase en
el modelo ideal y en el sector angular frontal, se utiliza la
condición

\begin{equation}
    d \leq \frac{\lambda_0}{2}.
    \label{eq:condicion_aliasing}
\end{equation}

Esta condición limita la diferencia de fase geométrica al intervalo
principal de fase. No elimina errores de calibración, saltos de ciclo,
reflexiones ni ambigüedades debidas a una relación señal-ruido
insuficiente.

Para una frecuencia de referencia
$f_0=\SI{2.66}{\kilo\hertz}$ y una velocidad del sonido de
$c=\SI{343.2}{\metre\per\second}$, se obtiene

\begin{equation}
    \lambda_0 \simeq \SI{12.9}{\centi\metre},
    \qquad
    \frac{\lambda_0}{2}\simeq \SI{6.45}{\centi\metre}.
\end{equation}

Una línea de base de $\SI{5.0}{\centi\metre}$ satisface la condición
anterior para esa frecuencia nominal. Esta elección reduce la
posibilidad de ambigüedad espacial idealizada en el sector frontal,
pero no constituye por sí sola una garantía de exactitud angular.

% ==============================================================================
% 2.4. ESTIMACIÓN DE VELOCIDAD MEDIANTE EL EFECTO DOPPLER
% ==============================================================================

\subsection{Estimación de velocidad mediante el efecto Doppler}

La velocidad de una fuente sonora puede estimarse a partir del
desplazamiento de frecuencia producido por su movimiento relativo
respecto a los receptores. Se considera una trayectoria rectilínea
entre dos receptores fijos $R_1$ y $R_2$, situados en $x=0$ y $x=L$,
respectivamente. La fuente se desplaza en el sentido positivo del
eje $x$, con velocidad $v>0$ cuando se aproxima a $R_2$ y se aleja
de $R_1$.

\begin{figure}[!htbp]
    \centering
    % figures/fig6_doppler_fdoa.tex
\begin{tikzpicture}[
    >=Stealth,
    sensor/.style={
        circle, draw=black!85, fill=black!12, line width=1pt,
        minimum size=0.55cm, inner sep=0pt
    },
    source/.style={
        circle, draw=black!85, fill=orange!80, line width=1.1pt,
        minimum size=0.55cm, inner sep=0pt
    },
    axis/.style={
        ->, line width=0.85pt, draw=black!65
    },
    velarrow/.style={
        ->, line width=1.3pt, draw=orange!90!black
    }
]

    % =========================================================================
    % 1. EJE 1D DE MOVIMIENTO Y RECEPTORES OPUESTOS R1 y R2 (y = 0)
    % =========================================================================
    % Eje horizontal
    \draw[axis] (-5.6, 0.0) -- (5.6, 0.0) node[right=2pt, font=\small\sffamily] {$x$ (Trayectoria 1D)};

    % Receptor R1 en x = -4.5 (Izquierda)
    \coordinate (R1) at (-4.5, 0.0);
    \node[sensor] (NodeR1) at (R1) {};
    \node[below=6pt, font=\small\sffamily\bfseries] at (R1) {$R_1$};
    \node[below=18pt, font=\footnotesize\sffamily, text=black!70] at (R1) {($x = 0$)};

    % Receptor R2 en x = +4.5 (Derecha)
    \coordinate (R2) at (4.5, 0.0);
    \node[sensor] (NodeR2) at (R2) {};
    \node[below=6pt, font=\small\sffamily\bfseries] at (R2) {$R_2$};
    \node[below=18pt, font=\footnotesize\sffamily, text=black!70] at (R2) {($x = L$)};

    % =========================================================================
    % 2. FUENTE EN MOVIMIENTO Y FLECHA DE VELOCIDAD
    % =========================================================================
    % Fuente S en x = 0.2
    \coordinate (S) at (0.2, 0.0);
    \node[source] (NodeS) at (S) {};
    \node[below=6pt, font=\small\sffamily\bfseries] at (S) {$S(t)$};

    % Flecha de velocidad instantánea v(t) elevada y acotada
    \draw[velarrow] (0.2, 0.55) -- ++(1.15, 0.0) node[above=2pt, font=\small\sffamily\bfseries] {$\vec{v}(t)$};

    % =========================================================================
    % 3. FRENTES DE ONDA DOPPLER EXCÉNTRICOS
    % =========================================================================
    % Centros de emisión pasados: x = -0.1, -0.4, -0.7
    % Radios: R = 0.95, 1.75, 2.55 (apertura controlada para no interferir)

    % --- Frentes dilatados hacia R1 (Izquierda: lambda1 > lambda0) ---
    % Cruzan el eje x en: -1.05, -2.15, -3.25
    \draw[line width=0.85pt, red!75!black] ($(-0.1, 0.0)+(158:0.95)$) arc (158:202:0.95);
    \draw[line width=0.85pt, red!75!black] ($(-0.4, 0.0)+(158:1.75)$) arc (158:202:1.75);
    \draw[line width=0.85pt, red!75!black] ($(-0.7, 0.0)+(158:2.55)$) arc (158:202:2.55);

    % --- Frentes comprimidos hacia R2 (Derecha: lambda2 < lambda0) ---
    % Cruzan el eje x en: +0.85, +1.35, +1.85
    \draw[line width=0.85pt, blue!75!black] ($(-0.1, 0.0)+(-22:0.95)$) arc (-22:22:0.95);
    \draw[line width=0.85pt, blue!75!black] ($(-0.4, 0.0)+(-22:1.75)$) arc (-22:22:1.75);
    \draw[line width=0.85pt, blue!75!black] ($(-0.7, 0.0)+(-22:2.55)$) arc (-22:22:2.55);

    % =========================================================================
    % 4. COTAS DE LONGITUD DE ONDA HOMOGÉNEAS (Ubicadas en y = -1.20)
    % =========================================================================
    % --- Cota Roja: Dilatación (lambda1 > lambda0) ---
    \draw[dotted, line width=0.5pt, draw=red!75!black] (-2.15, 0.0) -- (-2.15, -1.20);
    \draw[dotted, line width=0.5pt, draw=red!75!black] (-1.05, 0.0) -- (-1.05, -1.20);
    \draw[<->, draw=red!75!black, line width=0.65pt] (-2.15, -1.20) -- node[below=2pt, font=\footnotesize\sffamily] {$\lambda_1 > \lambda_0$} (-1.05, -1.20);

    % --- Cota Azul: Compresión (lambda2 < lambda0, idéntico formato) ---
    \draw[dotted, line width=0.5pt, draw=blue!75!black] (0.85, 0.0) -- (0.85, -1.20);
    \draw[dotted, line width=0.5pt, draw=blue!75!black] (1.35, 0.0) -- (1.35, -1.20);
    \draw[<->, draw=blue!75!black, line width=0.65pt] (0.85, -1.20) -- node[below=2pt, font=\footnotesize\sffamily] {$\lambda_2 < \lambda_0$} (1.35, -1.20);

    % =========================================================================
    % 5. CAJA SINTÉTICA DE CANCELACIÓN DE DERIVA (y = 3.65)
    % =========================================================================
    \node[draw=black!75, fill=white, rounded corners=4pt, line width=0.8pt, inner sep=7pt] at (0.0, 3.65) {
        \footnotesize
        \begin{tabular}{@{} l l @{}}
            \multicolumn{2}{@{} l @{}}{\textbf{Cinemática Doppler Diferencial en Régimen Subsónico ($v \ll c$):}} \\[4pt]
            Frecuencia en $R_1$ (Alejamiento): & $f_1(t) \approx f_0 \left(1 - \frac{v(t)}{c(T)}\right) \implies f_1 < f_0$ \\[2pt]
            Frecuencia en $R_2$ (Acercamiento): & $f_2(t) \approx f_0 \left(1 + \frac{v(t)}{c(T)}\right) \implies f_2 > f_0$ \\[5pt]
            \multicolumn{2}{@{} l @{}}{\textcolor{UdeAGreen!80!black}{\textbf{Operador diferencial covariante:}} \ $v(t) = c(T) \cdot \left(\frac{f_2(t) - f_1(t)}{f_1(t) + f_2(t)}\right) \quad \implies \text{Inmunidad a la deriva del emisor}$}
        \end{tabular}
    };

\end{tikzpicture}
    \vspace{0.25cm}
    \caption{Esquema del efecto Doppler para una fuente que se desplaza
    entre dos receptores fijos. La frecuencia recibida aumenta en el
    receptor hacia el que se aproxima la fuente y disminuye en el
    receptor del que se aleja.}
    \label{fig:doppler_fdoa}
\end{figure}

\subsubsection{Relación Doppler para receptores opuestos}

Para una fuente que emite una señal armónica de frecuencia $f_e$ y
se mueve con velocidad constante en un medio en reposo, las
frecuencias recibidas por los dos sensores son

\begin{equation}
    f_1 =
    \frac{f_e}{1+\dfrac{v}{c(T)}},
    \qquad
    f_2 =
    \frac{f_e}{1-\dfrac{v}{c(T)}}.
    \label{eq:doppler_exacto}
\end{equation}

Estas expresiones corresponden al caso ideal en que el movimiento
es rectilíneo, la velocidad es colineal con la línea fuente-receptor
y la frecuencia emitida puede considerarse constante durante los
intervalos de emisión que originan ambas señales recibidas. En un
movimiento acelerado, la formulación debe considerar las velocidades
y los instantes de emisión retardados correspondientes a cada
receptor.

Definiendo el número de Mach de la fuente como

\begin{equation}
    \beta = \frac{v}{c(T)},
\end{equation}

y considerando velocidades mucho menores que la velocidad del
sonido, $|\beta|\ll1$, las expresiones anteriores pueden aproximarse
a primer orden:

\begin{equation}
    f_1 \simeq f_e(1-\beta),
    \qquad
    f_2 \simeq f_e(1+\beta).
    \label{eq:doppler_unilateral}
\end{equation}

Esta aproximación describe el régimen de baja velocidad respecto
a la velocidad de propagación. Las correcciones de orden superior
no deben confundirse con la incertidumbre experimental del estimador.

\subsubsection{Estimador diferencial de velocidad}

A partir de las dos frecuencias recibidas, se define el estimador
diferencial

\begin{equation}
    \hat{v}_{\mathrm{diff}} =
    c(T)\frac{f_2-f_1}{f_1+f_2}.
    \label{eq:doppler_diferencial_formal}
\end{equation}

Si se sustituyen las expresiones exactas de la ecuación
\eqref{eq:doppler_exacto}, con una misma frecuencia emitida $f_e$
para ambos receptores, se obtiene

\begin{align}
    f_2-f_1
    &=
    f_e\left[
    \frac{1}{1-\beta}-\frac{1}{1+\beta}
    \right]
    =
    \frac{2f_e\beta}{1-\beta^2},
    \\
    f_1+f_2
    &=
    f_e\left[
    \frac{1}{1-\beta}+\frac{1}{1+\beta}
    \right]
    =
    \frac{2f_e}{1-\beta^2}.
\end{align}

Por tanto,

\begin{equation}
    \frac{f_2-f_1}{f_1+f_2}=\beta,
\end{equation}

y el estimador recupera

\begin{equation}
    \hat{v}_{\mathrm{diff}}=v.
    \label{eq:velocidad_final_doppler}
\end{equation}

La cancelación de $f_e$ es exacta dentro de este modelo. En particular,
el cociente no necesita conocer de antemano el valor absoluto de la
frecuencia emitida. Esta propiedad hace al estimador menos sensible
a cambios comunes de frecuencia que afecten por igual a las dos
mediciones.

Sin embargo, en una adquisición física los receptores pueden registrar
emisiones correspondientes a instantes distintos, debido a las
diferencias de distancia entre la fuente y cada sensor. Si la frecuencia
emitida varía con el tiempo, las dos frecuencias recibidas no contienen
necesariamente el mismo factor $f_e$. La cancelación deja entonces
de ser exacta.

Su validez como aproximación requiere que la deriva de frecuencia
sea suficientemente lenta durante la diferencia entre los tiempos
de emisión asociados a ambos receptores y que las ventanas de
estimación espectral sean temporalmente comparables.

La frecuencia media

\begin{equation}
    f_{\mathrm{mean}} =
    \frac{f_1+f_2}{2}
    \label{eq:f0_common}
\end{equation}

puede utilizarse como indicador de la escala de frecuencia común.
No coincide exactamente con la frecuencia emitida en el modelo
Doppler exacto, pues

\begin{equation}
    f_{\mathrm{mean}} =
    \frac{f_e}{1-\beta^2}.
\end{equation}

Para $|\beta|\ll1$, esta diferencia es de segundo orden y la
semisuma constituye una aproximación a la frecuencia de emisión.
Su uso como indicador de deriva debe interpretarse teniendo en
cuenta esta corrección y los posibles retardos de propagación
diferentes.

En una medición real, la superposición de trayectorias acústicas
puede afectar la frecuencia estimada. Si existen reflexiones,
cada trayectoria puede presentar una dirección de propagación
y un desplazamiento Doppler diferentes. Por tanto, la frecuencia
instantánea estimada no tiene por qué corresponder exclusivamente
a la trayectoria directa. La relación anterior debe entenderse
como el modelo de referencia que se contrasta experimentalmente,
no como una garantía de inmunidad frente a la propagación
multitrayecto.

% ==============================================================================
% 2.5. IDENTIFICABILIDAD Y RECONSTRUCCIÓN GEOMÉTRICA
% ==============================================================================

\subsection{Identificabilidad de la posición en dos dimensiones}

La estimación de una distancia, una dirección o una diferencia
de distancias no equivale necesariamente a determinar una posición
bidimensional. La reconstrucción espacial depende del número de
observables independientes, de la geometría de los receptores y
de las restricciones adicionales disponibles.

\subsubsection{Intersección de circunferencias a partir de ToA}

Dos distancias absolutas medidas desde receptores cuya posición
es conocida permiten reconstruir la posición de una fuente en
el plano, salvo una ambigüedad de simetría respecto a la línea
que une los receptores. Este procedimiento corresponde a la
intersección de dos circunferencias.

Se considera el sistema de referencia de la
Figura~\ref{fig:geometria_2d_angular}, con los receptores situados en

\begin{equation}
    R_1=\left(-\frac{d}{2},0\right),
    \qquad
    R_2=\left(+\frac{d}{2},0\right).
\end{equation}

Sean $r_1$ y $r_2$ las distancias estimadas desde la fuente hasta
los receptores $R_1$ y $R_2$. La posición $(x,y)$ satisface

\begin{align}
    \left(x+\frac{d}{2}\right)^2+y^2 &= r_1^2,
    \label{eq:circulo_r1}\\
    \left(x-\frac{d}{2}\right)^2+y^2 &= r_2^2.
    \label{eq:circulo_r2}
\end{align}

Restando ambas ecuaciones se obtiene

\begin{equation}
    x=\frac{r_1^2-r_2^2}{2d}.
    \label{eq:coordenada_x_analitica}
\end{equation}

La coordenada $y$ se obtiene sustituyendo este resultado en cualquiera
de las ecuaciones de circunferencia:

\begin{equation}
    y =
    \pm\sqrt{
    r_1^2-\left(x+\frac{d}{2}\right)^2
    }.
    \label{eq:coordenada_y_analitica}
\end{equation}

El signo doble representa las dos posiciones simétricas que producen
las mismas distancias a los receptores. La selección de la solución
con $y>0$ requiere información adicional sobre el semiplano en el
que se encuentra la fuente; no puede deducirse únicamente de las
dos distancias.

Con mediciones ruidosas, las circunferencias pueden no intersectarse
exactamente. En ese caso, el radicando de la ecuación
\eqref{eq:coordenada_y_analitica} puede resultar negativo o muy
próximo a cero. No es apropiado corregir automáticamente ese
resultado mediante un recorte numérico sin evaluar antes la
compatibilidad de las medidas. Una reconstrucción robusta debe
considerar las incertidumbres de ambas distancias y, si es
necesario, resolver un problema de mínimos cuadrados.

\subsubsection{Diferencia de tiempos de llegada y geometría hiperbólica}

Cuando el observable disponible es la diferencia de tiempos de
llegada entre dos receptores, puede calcularse una diferencia
de distancias

\begin{equation}
    \Delta r_{\mathrm{TDOA}} =
    c(T)(t_1-t_2).
    \label{eq:diferencia_distancias_tdoa}
\end{equation}

El conjunto de puntos que satisface

\begin{equation}
    r_1-r_2=\Delta r_{\mathrm{TDOA}}
    \label{eq:hiperbola_tdoa}
\end{equation}

forma una hipérbola cuyos focos están en las posiciones de los
dos receptores. Una única diferencia de tiempos no determina
por sí sola una posición bidimensional única: la fuente puede
encontrarse en cualquier punto de esa hipérbola.

Se requiere información adicional, como una segunda diferencia
independiente, una distancia absoluta o una restricción geométrica
sobre la trayectoria. Por consiguiente, la reconstrucción por
intersección de circunferencias y la localización hiperbólica
TDOA son formulaciones relacionadas, pero no intercambiables.
La primera utiliza dos distancias absolutas; la segunda utiliza
diferencias de distancias.

\subsubsection{Ambigüedad de fase y diferencia de camino}

La diferencia de fase de una señal armónica proporciona una
diferencia de camino módulo la longitud de onda:

\begin{equation}
    \Delta r_{\Phi} =
    \frac{\lambda_0}{2\pi}
    \left(\Delta\Phi+2\pi n\right),
    \qquad n\in\mathbb{Z}.
    \label{eq:ambiguedad_fase}
\end{equation}

El entero $n$ representa la ambigüedad de ciclo. La fase por sí
sola no determina dicho entero. Para resolverlo se requiere
información adicional, por ejemplo una estimación gruesa obtenida
mediante ToA, una restricción sobre el desplazamiento entre
mediciones sucesivas o una condición geométrica que limite la
diferencia de camino.

La operación de reducción módulo $\lambda_0$ puede representar
la fase principal, pero no recupera por sí misma la diferencia
de camino absoluta ni garantiza continuidad temporal. En
particular, no debe aplicarse una corrección modular a una
diferencia de distancias ya obtenida mediante ToA, salvo que
se esté resolviendo explícitamente una ambigüedad de fase y
se compruebe la compatibilidad con las distancias absolutas.

\subsubsection{Alcance geométrico del arreglo utilizado}

Las relaciones anteriores describen modelos geométricos generales.
Sin embargo, la geometría del instrumento impone límites a las
magnitudes que pueden recuperarse.

Con dos receptores dispuestos sobre una única línea base, el
TDOA define una diferencia de distancias y el PDoA permite
estimar una dirección dentro del plano de incidencia considerado.
Ninguno de estos observables, por separado, determina una
posición bidimensional única.

La combinación de una distancia radial y una dirección puede
proporcionar una posición en un sistema de coordenadas definido,
siempre que ambas estimaciones correspondan al mismo estado de
la fuente y que sus referencias geométricas sean compatibles.
Esta combinación no debe confundirse con una reconstrucción
multilaterada obtenida a partir de varias líneas base
independientes.

Por tanto, el alcance espacial debe establecerse a partir de
los observables efectivamente disponibles y de la geometría
del montaje, sin atribuir al arreglo una capacidad de
localización que no se desprenda de esas condiciones.

% ==============================================================================
% 2.6. COMPARACIÓN Y CONDICIONES DE VALIDEZ
% ==============================================================================

\subsection{Comparación de observables y condiciones de validez}

Los métodos descritos utilizan magnitudes físicas diferentes y
presentan sensibilidades distintas frente a las condiciones de
propagación, la geometría y la adquisición. La
Tabla~\ref{tab:comparacion_observables} resume sus principales
características.

\begin{table}[!htbp]
    \centering
    \small
    \renewcommand{\arraystretch}{1.25}
    \begin{tabularx}{\textwidth}{
        >{\raggedright\arraybackslash}p{1.8cm}
        >{\raggedright\arraybackslash}p{2.0cm}
        X
        X
    }
        \toprule
        \textbf{Método} &
        \textbf{Observable} &
        \textbf{Información obtenida} &
        \textbf{Limitación dominante}
        \\
        \midrule

        ToA &
        $t_{\mathrm{arr}}$ &
        Distancia absoluta cuando se conoce el instante de emisión y
        se corrigen los retardos instrumentales. &
        Sensibilidad a la detección del arribo, la sincronización y
        la calibración temporal.
        \\

        TDOA &
        $t_1-t_2$ &
        Diferencia de distancias entre la fuente y los receptores;
        permite estimar dirección bajo una geometría conocida. &
        Una sola diferencia define un lugar geométrico y no una
        posición bidimensional única.
        \\

        Amplitud y energía &
        $A(r)$, $\mathcal{E}_{\mathrm{dig}}(r)$ &
        Estimación radial mediante modelos empíricos calibrados. &
        Sensibilidad a la directividad, la ganancia, las reflexiones
        y la definición de la ventana temporal.
        \\

        PDoA &
        $\Delta\Phi$ &
        Dirección de llegada a partir de la diferencia de fase
        entre receptores. &
        Ambigüedad de fase, calibración intercanal y propagación
        multitrayecto.
        \\

        Doppler diferencial &
        $f_1$, $f_2$ &
        Velocidad relativa bajo una geometría de receptores
        opuestos y frecuencias comparables. &
        Requiere geometría conocida, estimaciones espectrales
        comparables y deriva suficientemente lenta.
        \\

        \bottomrule
    \end{tabularx}
    \caption{Comparación cualitativa de los observables utilizados
    para estimar distancia, dirección y velocidad. Las prestaciones
    cuantitativas deben establecerse mediante la caracterización
    experimental de cada estimador.}
    \label{tab:comparacion_observables}
\end{table}

La combinación de estos observables permite abordar la localización
acústica desde perspectivas complementarias. Los tiempos de llegada
aportan información sobre la distancia recorrida; las diferencias
de tiempo y fase relacionan las señales registradas por sensores
separados; las magnitudes de amplitud y energía permiten construir
estimadores radiales calibrados; y las frecuencias recibidas
contienen información sobre el movimiento relativo entre la fuente
y los receptores.

La calidad de una reconstrucción conjunta depende de la consistencia
entre estos observables, de la geometría del montaje y de la
caracterización de sus incertidumbres. En consecuencia, las
relaciones analíticas de este capítulo establecen los modelos de
inversión que deben contrastarse con los resultados experimentales.

% ==============================================================================
% FIN DEL CAPÍTULO 2
% ==============================================================================
